\documentclass[11pt]{article}
% Shared preamble for manuscripts-v2 papers. Load with % Shared preamble for manuscripts-v2 papers. Load with % Shared preamble for manuscripts-v2 papers. Load with \input{../common/preamble.tex}
% immediately after \documentclass[11pt]{article}.
\ifdefined\pdftrailerid\pdftrailerid{}\fi
\usepackage[T1]{fontenc}
\usepackage{lmodern}
\usepackage[letterpaper,margin=1.02in]{geometry}
\usepackage{amsmath,amssymb,amsthm,mathtools,booktabs,array,microtype,enumitem}
\usepackage{tikz}
\usepackage{xurl}
\usepackage[colorlinks=true,allcolors=bluegreen]{hyperref}
\definecolor{bluegreen}{RGB}{22,83,91}
\usepackage[capitalize,nameinlink]{cleveref}

\theoremstyle{plain}
\newtheorem{theorem}{Theorem}[section]
\newtheorem{lemma}[theorem]{Lemma}
\newtheorem{proposition}[theorem]{Proposition}
\newtheorem{corollary}[theorem]{Corollary}
\newtheorem{conjecture}[theorem]{Conjecture}
\newtheorem*{theorem*}{Theorem}
\newtheorem{maintheorem}{Theorem}
\renewcommand{\themaintheorem}{\Alph{maintheorem}}
\theoremstyle{definition}
\newtheorem{definition}[theorem]{Definition}
\newtheorem{example}[theorem]{Example}
\newtheorem{question}[theorem]{Question}
\theoremstyle{remark}
\newtheorem{remark}[theorem]{Remark}
\crefname{maintheorem}{Theorem}{Theorems}
\crefname{conjecture}{Conjecture}{Conjectures}
\crefname{question}{Question}{Questions}
\crefname{example}{Example}{Examples}

\newcommand{\R}{\mathbb R}
\newcommand{\Q}{\mathbb Q}
\newcommand{\Z}{\mathbb Z}
\newcommand{\N}{\mathbb N}
\newcommand{\T}{\mathbb T}
\newcommand{\C}{\mathbb C}
\DeclareMathOperator{\den}{den}
\newcommand{\e}{\mathrm e}

% Immutable Lean source links at research-checkpoint-18.
\newcommand{\leancommit}{e4809cf561573aef588e6ffc491a35cdb79fd9ec}
\newcommand{\leanbase}{https://github.com/eldernos/MathematicsResearch/blob/\leancommit}
% \lean{File}{line}{label}: link to NormalityChecker/File.lean#Lline at the pinned commit.
\newcommand{\lean}[3]{\href{\leanbase/NormalityChecker/#1.lean\#L#2}{#3}}
% \repo{path}{label}: link to any repository path at the pinned commit.
\newcommand{\repo}[2]{\href{https://github.com/eldernos/MathematicsResearch/tree/\leancommit/#1}{#2}}

% Evidence labels (see papers/WAVE2-CONTRACT.md).
\newcommand{\evidence}[1]{\textup{[\emph{#1}]}}

\setlength{\emergencystretch}{1.5em}
\raggedbottom

% immediately after \documentclass[11pt]{article}.
\ifdefined\pdftrailerid\pdftrailerid{}\fi
\usepackage[T1]{fontenc}
\usepackage{lmodern}
\usepackage[letterpaper,margin=1.02in]{geometry}
\usepackage{amsmath,amssymb,amsthm,mathtools,booktabs,array,microtype,enumitem}
\usepackage{tikz}
\usepackage{xurl}
\usepackage[colorlinks=true,allcolors=bluegreen]{hyperref}
\definecolor{bluegreen}{RGB}{22,83,91}
\usepackage[capitalize,nameinlink]{cleveref}

\theoremstyle{plain}
\newtheorem{theorem}{Theorem}[section]
\newtheorem{lemma}[theorem]{Lemma}
\newtheorem{proposition}[theorem]{Proposition}
\newtheorem{corollary}[theorem]{Corollary}
\newtheorem{conjecture}[theorem]{Conjecture}
\newtheorem*{theorem*}{Theorem}
\newtheorem{maintheorem}{Theorem}
\renewcommand{\themaintheorem}{\Alph{maintheorem}}
\theoremstyle{definition}
\newtheorem{definition}[theorem]{Definition}
\newtheorem{example}[theorem]{Example}
\newtheorem{question}[theorem]{Question}
\theoremstyle{remark}
\newtheorem{remark}[theorem]{Remark}
\crefname{maintheorem}{Theorem}{Theorems}
\crefname{conjecture}{Conjecture}{Conjectures}
\crefname{question}{Question}{Questions}
\crefname{example}{Example}{Examples}

\newcommand{\R}{\mathbb R}
\newcommand{\Q}{\mathbb Q}
\newcommand{\Z}{\mathbb Z}
\newcommand{\N}{\mathbb N}
\newcommand{\T}{\mathbb T}
\newcommand{\C}{\mathbb C}
\DeclareMathOperator{\den}{den}
\newcommand{\e}{\mathrm e}

% Immutable Lean source links at research-checkpoint-18.
\newcommand{\leancommit}{e4809cf561573aef588e6ffc491a35cdb79fd9ec}
\newcommand{\leanbase}{https://github.com/eldernos/MathematicsResearch/blob/\leancommit}
% \lean{File}{line}{label}: link to NormalityChecker/File.lean#Lline at the pinned commit.
\newcommand{\lean}[3]{\href{\leanbase/NormalityChecker/#1.lean\#L#2}{#3}}
% \repo{path}{label}: link to any repository path at the pinned commit.
\newcommand{\repo}[2]{\href{https://github.com/eldernos/MathematicsResearch/tree/\leancommit/#1}{#2}}

% Evidence labels (see papers/WAVE2-CONTRACT.md).
\newcommand{\evidence}[1]{\textup{[\emph{#1}]}}

\setlength{\emergencystretch}{1.5em}
\raggedbottom

% immediately after \documentclass[11pt]{article}.
\ifdefined\pdftrailerid\pdftrailerid{}\fi
\usepackage[T1]{fontenc}
\usepackage{lmodern}
\usepackage[letterpaper,margin=1.02in]{geometry}
\usepackage{amsmath,amssymb,amsthm,mathtools,booktabs,array,microtype,enumitem}
\usepackage{tikz}
\usepackage{xurl}
\usepackage[colorlinks=true,allcolors=bluegreen]{hyperref}
\definecolor{bluegreen}{RGB}{22,83,91}
\usepackage[capitalize,nameinlink]{cleveref}

\theoremstyle{plain}
\newtheorem{theorem}{Theorem}[section]
\newtheorem{lemma}[theorem]{Lemma}
\newtheorem{proposition}[theorem]{Proposition}
\newtheorem{corollary}[theorem]{Corollary}
\newtheorem{conjecture}[theorem]{Conjecture}
\newtheorem*{theorem*}{Theorem}
\newtheorem{maintheorem}{Theorem}
\renewcommand{\themaintheorem}{\Alph{maintheorem}}
\theoremstyle{definition}
\newtheorem{definition}[theorem]{Definition}
\newtheorem{example}[theorem]{Example}
\newtheorem{question}[theorem]{Question}
\theoremstyle{remark}
\newtheorem{remark}[theorem]{Remark}
\crefname{maintheorem}{Theorem}{Theorems}
\crefname{conjecture}{Conjecture}{Conjectures}
\crefname{question}{Question}{Questions}
\crefname{example}{Example}{Examples}

\newcommand{\R}{\mathbb R}
\newcommand{\Q}{\mathbb Q}
\newcommand{\Z}{\mathbb Z}
\newcommand{\N}{\mathbb N}
\newcommand{\T}{\mathbb T}
\newcommand{\C}{\mathbb C}
\DeclareMathOperator{\den}{den}
\newcommand{\e}{\mathrm e}

% Immutable Lean source links at research-checkpoint-18.
\newcommand{\leancommit}{e4809cf561573aef588e6ffc491a35cdb79fd9ec}
\newcommand{\leanbase}{https://github.com/eldernos/MathematicsResearch/blob/\leancommit}
% \lean{File}{line}{label}: link to NormalityChecker/File.lean#Lline at the pinned commit.
\newcommand{\lean}[3]{\href{\leanbase/NormalityChecker/#1.lean\#L#2}{#3}}
% \repo{path}{label}: link to any repository path at the pinned commit.
\newcommand{\repo}[2]{\href{https://github.com/eldernos/MathematicsResearch/tree/\leancommit/#1}{#2}}

% Evidence labels (see papers/WAVE2-CONTRACT.md).
\newcommand{\evidence}[1]{\textup{[\emph{#1}]}}

\setlength{\emergencystretch}{1.5em}
\raggedbottom

% Author information approved by the repository owner on 2026-11-25:
%   name: Christopher Anderson; designation: Independent Researcher.
% No academic affiliation, ORCID or personal contact details have been supplied.
% Do not add any that the owner has not approved.
\newcommand{\OwnerAuthor}{Christopher Anderson}
\newcommand{\OwnerAffiliation}{Independent Researcher}
\newcommand{\OwnerContact}{Corrections and comments: \url{https://github.com/eldernos/MathematicsResearch/issues}}
\newcommand{\ReleaseLabel}{manuscripts-v2.0.0}
\newcommand{\LicenseLine}{\copyright\ 2026 Christopher Anderson. Licensed under \href{https://creativecommons.org/licenses/by/4.0/}{CC BY 4.0}.}
% \makebyline{version}{date}: centred author block used under the title.
\newcommand{\makebyline}[2]{%
  {\large \OwnerAuthor\par}%
  \ifx\OwnerAffiliation\empty\else{\small \OwnerAffiliation\par}\fi
  \vspace{0.3em}
  {\small Version #1 \enspace$\cdot$\enspace #2 \enspace$\cdot$\enspace \ReleaseLabel\par}%
  {\footnotesize \OwnerContact\par}%
  {\footnotesize \LicenseLine\par}%
  \hypersetup{pdfauthor={\OwnerAuthor}}}

\newcommand{\circabs}[1]{\lVert #1\rVert_{\T}}
\hypersetup{pdftitle={Arithmetic Certificates for Normality and Mathematical Limitations}}

\begin{document}
\begin{center}
{\LARGE\bfseries Arithmetic Certificates for Normality\\and Mathematical Limitations\par}
\vspace{0.6em}
\makebyline{2.0}{25 November 2026}
\end{center}

\begin{abstract}
A radix orbit in base \(b\ge2\) is equidistributed if, for every positive tolerance and denominator floor, almost all blocks on a prescribed geometric schedule admit rational models with mean circle error below the tolerance, reduced denominators above the floor on a fixed prime support coprime to \(b\), and a polynomial denominator cap. Partial coverage gives a quantitative discrepancy bound. Accurate fixed-support models on long overlaps must agree modulo one after alignment, and a counting estimate (Theorem~\ref{thm:count}) shows that polynomial-denominator rational models meeting the mean circle error threshold \(1/(16b^4)\) on an independently prescribed geometric schedule with polynomially many starts per scale occur only finitely often for almost every real. Explicit pairs of binary streams separate odd-primary rational data and logarithmic word laws from ordinary normality. Two uniform halting reductions show that neither normality nor nonnormality is computably enumerable on digit-program names, even when their values are transcendental and their binary discrepancy tends to zero or one. Lean~4 verifies the underlying fixed-support estimate, uniform overlap separation, the certificate transfer step, radix-orbit partial-coverage bounds, the finite weighted word inequality used in the counting proof, and---under explicit endpoint hypotheses---each individual finite-edit, halting-reduction, and nonrecognizability declaration. A separate written Rosser argument establishes consistency-relative independence for the normality sentence of a theory-dependent digit program, establishing neither normality nor independence for any named classical constant.
\end{abstract}

\section{Normality from rational orbit models}
Let \(b\ge2\) be an integer and let \(x\in\R\). Write \(x_n=\{b^n x\}\in[0,1)\) for integers \(n\ge0\), \(e(t)=\exp(2\pi i t)\), and
\(\circabs{t}=\min_{m\in\Z}|t-m|\). For every integer \(N\ge1\) set
\[
 D_N^{(b)}(x)=\sup_{0\le u<v\le1}
 \left|\frac1N\#\{0\le n<N:x_n\in[u,v)\}-(v-u)\right|,
 \qquad D_\infty^{(b)}(x)=\limsup_{N\to\infty}D_N^{(b)}(x).
\]
A real number is normal in base \(b\) if every word of length \(r\ge1\) has limiting frequency \(b^{-r}\), counting overlapping occurrences and using the terminating expansion at radix rationals. Each word corresponds to a half-open interval \([m/b^r,(m+1)/b^r)\), where \(0\le m<b^r\). Approximating both endpoints of any interval by points of this grid shows that normality is equivalent to \(D_N^{(b)}(x)\to0\). Write \(\den y\) for the positive denominator of \(y\in\Q\) in lowest terms. Call an integer \(Q\) supported on \(c\) if \(Q\mid c^j\) for some integer \(j\ge0\).

\begin{theorem}[Mean rational approximation]\label{thm:mean}
\evidence{Lean-supported paper proof}
Fix integers \(b,c\ge2\) with \(\gcd(b,c)=1\), and integers \(L\ge2,\ B\ge1,\ d\ge0\). Let \(x\in\R\) and let \(P_n\in\Q\) satisfy \(b^nx-P_n\to0\) in \(\R\). At integer scale \(k\ge0\) put
\[
 A=(LB)^k,\qquad T=B^k,\qquad m=L^k,\qquad A=mT.
\]
For each integer \(u\in\{1,\ldots,LB-1\}\), integer \(M\ge1\) and real \(\eta>0\), suppose that all sufficiently large integers \(k\) admit a set \(E\subseteq\{0,\ldots,m-1\}\) with \(|E|\le\eta m\) such that for every integer \(j\in\{0,\ldots,m-1\}\setminus E\), writing \(n=uA+jT\), there exists a rational \(y\) with reduced denominator \(Q\) supported on \(c\) satisfying
\begin{equation}\label{eq:mean}
 M\le Q\le T^d,\qquad
 \frac1T\sum_{i=0}^{T-1}\circabs{P_{n+i}-b^i y}\le\eta .
\end{equation}
Then \(x\) is normal in base \(b\).
\end{theorem}

An \emph{arithmetic certificate} in this setting consists of the block models, the exceptional sets and the estimates in \eqref{eq:mean}, for fixed \(b,c,L,B,d\) and every \(M,\eta\). In the exact case \(y=P_n\) and \(P_{n+i}=b^iP_n\) throughout each retained block. If \(B=1\) or \(d=0\), then \(T^d=1\) precludes the hypotheses for \(M\ge2\) and \(\eta<1\). The real tail transfers character-sum bounds from \(P_n\) to the radix shifts.

Korobov studies word frequencies in complete and incomplete periods of rational fractions \cite[pp.~659--660]{korobov}. Bailey and Crandall transfer equidistribution between perturbed radix recurrences and digit shifts, and apply this transfer to lacunary series \cite[Theorems~3.2 and 4.8]{bailey}. Wagner uses arithmetic constructions to produce rings with contrasting normality in related bases \cite[p.~211]{wagner}. Schmidt's theorem for multiplicatively independent bases uses exponential-sum estimates against a measure carried by a restricted digit construction \cite[Theorem~1]{schmidt}. Vandehey's series criterion uses a growing divisibility chain and an insertion-gap hypothesis \cite[Theorem~7.4, p.~24 of the preprint]{vandehey}. Here the denominators belong to reduced models of shifted prefixes and the error is averaged within a block. Related worked examples use sparse Stoneham series and localized logarithms. A localized-logarithm manuscript hosted by Captain Sude gives a quantitative example \cite[Theorems~1.1 and 4.1]{localized}; Mira Cult states analogous fixed-support estimates and reduced-prefix proofs for several sparse series \cite[Theorems~1--3]{mira}.

\subsection{Short sums and transfer to every prefix}
\begin{theorem}[Fixed-support estimate, qualitative form of Vandehey]\label{thm:short}
\evidence{Kernel-verified}
For coprime integers \(b,c\ge2\), each integer \(d\ge0\) and each real \(\epsilon>0\), there is an integer \(M=M(b,c,d,\epsilon)\ge1\) such that for every integer \(T\ge1\) and every \(y=a/Q\in\Q\) in lowest terms with \(Q\) supported on \(c\) and \(M\le Q\le T^d\),
\[
 \left|\frac1T\sum_{i=0}^{T-1}e(b^iy)\right|\le\epsilon .
\]
\end{theorem}
For \(d>0\) the range is \(T\ge Q^{1/d}\). Vandehey's Corollary~1.2 \cite[p.~2 of the preprint]{vandehey} gives \(O_{b,c,d}(Q^{-\sigma}(1+\log Q))\) for the displayed mean, for some \(\sigma>0\). His sum starts at exponent one, an equivalent convention since \(b\) permutes the reduced numerators. The following differencing argument proves this qualitative consequence directly.

\begin{proof}
Let \(r=\varphi(c^2)\) and \(C=b^r-1\). For each prime \(p\mid c\) and integer \(m\ge1\),
\[
 v_p(b^{rm}-1)=v_p(C)+v_p(m).
\]
Indeed \(p\mid C\) by Euler's theorem, and \(4\mid C\) if \(2\mid c\). For \(p\mid z\) (and \(4\mid z\) at \(2\)), the binomial expansion of \((1+z)^p-1\) gives valuation \(v_p(z)+1\). For \(p\nmid u\), the same expansion gives \(v_p((1+z)^u-1)=v_p(z)\). Iteration proves the formula. For \(Q>1\), let \(v_Q\) denote the order of \(b\) modulo \(Q\). Then \(Q\mid b^{rv_Q}-1\), hence \(Q\mid C v_Q\). In particular \(Q/C\le v_Q\le Q\).

Let \(G\) be the subgroup generated by \(b\) modulo \(Q\). Its Fourier sums satisfy
\[
 S(a)=\sum_{t\in G}e(at/Q),\qquad \sum_{a\bmod Q}|S(a)|^2=Qv_Q
\]
by character orthogonality. For a unit \(a\), \(S\) is constant on the \(v_Q\) points in \(aG\), so \(|S(a)|\le\sqrt Q\). Splitting into complete periods and a remainder yields
\begin{equation}\label{eq:period}
 \frac1T\left|\sum_{i<T}e(b^iy)\right|\le \frac{C}{\sqrt Q}+\frac{Q}{T}.
\end{equation}
Here and below \(i<T\) means \(0\le i<T\).

We prove by induction on the integer \(s\ge2\) that for every \(\epsilon>0\) there exist integers \(K_0\ge1\) and \(M\ge1\) such that the normalized sum has modulus at most \(\epsilon\) whenever \(K\ge K_0\), \(T\ge K^3\), and \(M\le Q\le K^s\) (for integers \(K,T,Q\)). Equation~\eqref{eq:period} proves the case \(s=2\), taking \(K_0\ge 2/\epsilon\) and \(M\ge (2C/\epsilon)^2\). Suppose the assertion holds for \(s\), and fix \(\epsilon>0\). First choose an integer \(H\ge1\) with \(1/H\le\epsilon^2/8\). Next set the induction tolerance \(\delta=\epsilon^2/8\), and invoke the induction hypothesis at \(s\) with tolerance \(\delta\) to obtain \(K_1,M_1\) and with tolerance \(\epsilon\) to obtain \(K_2,M_2\). Choose an integer \(K_0\ge\max(K_1,K_2)\) large enough that \(K_0^{s-1}/(CHc)\ge M_1\) and \(2Hc/K_0^2\le\epsilon/2\), and put \(M=M_2\).

Let \(K\ge K_0\), \(T\ge K^3\), and \(M\le Q\le K^{s+1}\). Values with \(Q\le K^s\) are already bounded by \(\epsilon\) from \((K_2,M_2)\), so assume \(K^s<Q\le K^{s+1}\). Multiply successive prime factors of \(Q\), with multiplicity, until their product \(D\) first reaches \(K\). Then \(D\mid Q\) and \(K\le D\le cK\). Write \(v=v_D\le cK\). For integers \(0\le j<l<H\) the difference of the translated phases \((b^{lv}-b^{jv})y\) has reduced denominator \(Q/g\), where
\[
 g=\gcd(Q,b^{(l-j)v}-1),\qquad D\mid g,\qquad g\le C(l-j)v\le CHcK.
\]
The last bound follows since \(v_g\mid(l-j)v\) and \(g\le C v_g\). Consequently every pair \(0\le j<l<H\) satisfies
\[
 M_1\le \frac{K^{s-1}}{CHc}<\frac{Q}{g}\le K^s.
\]
Set \(z_i=e(b^iy)\). The mean of \(z_{i+jv}\) over \(i<T\) differs from that of \(z_i\) by at most \(2Hv/T\le 2Hc/K^2\le\epsilon/2\). Average over \(0\le j<H\) and apply Cauchy--Schwarz in \(i\). The squared modulus of the translated mean is bounded by
\[
 \frac1H+\frac2{H^2}\sum_{0\le j<l<H}
       \left|\frac1T\sum_{i<T} z_{i+lv}\overline{z_{i+jv}}\right| .
\]
Because \(K\ge K_1\) and \(M_1\le Q/g\le K^s\), the induction hypothesis at tolerance \(\delta\) uniformly bounds every inner mean by \(\delta\). The squared modulus of the translated mean is therefore at most \(1/H+\delta\le\epsilon^2/4\), so the translated mean has modulus at most \(\epsilon/2\). Adding the translation error gives modulus at most \(\epsilon\), completing the induction step.

Finally, the case \(d=0\) is vacuous with \(M=2\). For \(d\ge1\), apply the induction statement with \(s=4d\) to obtain \(K_0,M_0\), and set \(K_1=\max(K_0,4)\) and \(M=\max(M_0,K_1^{3d})\). Whenever \(M\le Q\le T^d\), we have \(T\ge Q^{1/d}\ge K_1^3\), so \(K=\lfloor T^{1/3}\rfloor\) satisfies \(K\ge K_1\ge K_0\) and \(T\ge K^3\); the finitely many integers \(T<K_1^3\) are excluded by \(M\). Moreover \(K\ge4\) implies \(T<(K+1)^3\le K^4\), hence \(Q\le T^d\le K^{4d}\). The induction statement at \(s=4d\) then bounds the normalized sum by \(\epsilon\).
\end{proof}

\begin{proof}[Proof of Theorem~\ref{thm:mean}]
Fix a nonzero integer \(h\). The denominator of \(hy\) is \(Q/\gcd(Q,|h|)\), between \(Q/|h|\) and \(Q\). Apply Theorem~\ref{thm:short} with tolerance \(\eta\) and choose the floor in \eqref{eq:mean} at least \(|h|\) times its threshold. For any \(s,t\in\R\),
\[
 |e(hs)-e(ht)|\le 2\pi |h|\,\circabs{s-t}.
\]
Thus a good full block contributes at most \((1+2\pi |h|)\eta T\) to the modulus of the \(h\)-sum for \(P\). More generally, if model sums are at most \(\epsilon T\), mean errors are at most \(\delta\), and \(E\) is the exceptional set of block indices, then at every integer cut \(0\le t\le mT\)
\begin{equation}\label{eq:prefix}
 \left|\sum_{i<t}e(h P_{uA+i})\right|
 \le(\epsilon+2\pi |h|\delta)mT+(|E|+1)T .
\end{equation}
In fact, each exceptional full block and the unfinished last block contributes at most \(T\) to the absolute sum.

Since \(T/A=L^{-k}\to0\), \eqref{eq:prefix} bounds every partial cell sum by \(o(A)\), uniformly in the finitely many \(u\), on letting \(\eta\downarrow0\). For a prefix ending in the shell \([A,(LB)A)\), sum over its earlier shells and its completed or partial cells. From any fixed large scale onward, an upper bound \(\rho (LB)^j\) per cell contributes at most \(\rho (LB-1)\sum_{j=0}^k (LB)^j\le \rho LB\,A\). Earlier positions contribute a bounded term. Divide by the prefix length, at least \(A\), and send \(\rho\) to zero. The character means for \(P_n\) therefore vanish. Those for \(b^nx\) have the same limit by the tail hypothesis and the Lipschitz bound. Weyl's criterion \cite[\S1]{weyl} proves normality.
\end{proof}

\section{Partial coverage}\label{sec:cover}
Let \((x_i)_{i\ge0}\) be any sequence of observations in \([0,1)\). For any finite set \(F\subset\Z_{\ge0}\) of positions and any half-open test interval \(I\subset[0,1)\), write
\[
 C_I(F)=\sum_{i\in F}(\mathbf1_I(x_i)-|I|),
\]
with the empty-sum convention \(C_I(\varnothing)=0\). For a nonempty finite set \(F\) put
\[
 D(F)=\frac{1}{|F|}\sup_{0\le u<v\le1} |C_{[u,v)}(F)|,
\]
and for each integer \(N\ge1\) denote the prefix discrepancy and upper limiting discrepancy of \((x_i)\) by
\[
 D_N=D(\{0,\ldots,N-1\}),\qquad D_\infty=\limsup_{N\to\infty}D_N.
\]
The following estimates apply to any observations \(x_i\in[0,1)\) and need no arithmetic hypothesis on the pieces.

\begin{proposition}[Partial coverage]\label{prop:cover}
\evidence{Lean-supported paper proof}
Suppose disjoint pieces in \(\{0,\ldots,N-1\}\) cover an integer count of \(G\) positions altogether and each nonempty piece has discrepancy at most \(\delta\in[0,1]\). Then
\begin{equation}\label{eq:covered}
 D_N\le 1-(1-\delta)G/N .
\end{equation}
If, for every \(\epsilon>0\) and all sufficiently large integers \(N\), such pieces can be chosen with \(G/N\ge g-\epsilon\), then \(D_\infty\le1-g(1-\delta)\). More generally let \(F_N\subseteq\{0,\ldots,N-1\}\) with \(\liminf_{N\to\infty} |F_N|/N\ge g\). If \(\sup_I|C_I(F_N)|=o(N)\), then \(D_\infty\le1-g\).
\end{proposition}
\begin{proof}
On a piece \(F\), the centered sum has absolute value at most \(\delta |F|\) (trivially if \(F=\varnothing\)). An uncovered position contributes at most one in absolute value since \(|\mathbf1_I-|I||\le1\). Add over the disjoint pieces and the remainder, then divide by \(N\) and take the supremum. The asymptotic inequality follows on letting \(\epsilon\downarrow0\). In the last assertion the two analogous costs are \(o(N)\) and \(N-|F_N|\).
\end{proof}
A fixed digit resolution has a different effect. If each of the \(q\ge1\) equal mesh intervals \([j/q,(j+1)/q)\) has limiting frequency \(1/q\), every union of consecutive mesh intervals has its correct limiting frequency, uniformly over this finite list. For each test interval \(I\subset[0,1)\), the circumscribed union of mesh intervals exceeds \(|I|\) by at most \(2/q\), and the inscribed union of mesh intervals (which may be empty, with count and length zero) falls short of \(|I|\) by at most \(2/q\). Hence \(D_\infty\le2/q\). In particular a test at one resolution gives a discrepancy bound even when finer resolutions have yet to be controlled.

Proposition~\ref{prop:cover} is a weighted concatenation bound charging uncovered positions at one. For the full unsigned localized-support classification built on these coverage estimates, see the companion manuscript \emph{Normality of localized logarithms and polylogarithms on finite prime supports} \cite{companionlocalized}; pointers to the archived version~1.1 endpoint-path and \(\{2,5,13\}\) material are recorded in Section~9.

\section{Separation on a shared segment}\label{sec:rigidity}
\begin{theorem}[Uniform separation]\label{thm:rigidity}
\evidence{Kernel-verified}
Fix coprime integers \(b,c\ge2\) and an integer \(d\ge0\). There exists \(\gamma=\gamma(b,c,d)>0\) with the following property. Suppose \(T\ge1\) is an integer, \(y,z\in\Q\) have reduced denominators supported on \(c\) and at most \(T^d\), and \(y-z\notin\Z\). For any real observations \(f_0,\ldots,f_{T-1}\),
\begin{equation}\label{eq:rigid}
 \frac1T\sum_{i<T}\big(\circabs{f_i-b^iy}+\circabs{f_i-b^iz}\big)\ge\gamma .
\end{equation}
\end{theorem}
\begin{proof}
Put \(w=y-z\) and \(Q=\den w>1\). Then \(Q\) divides \(\den(y)\den(z)\), so its support is fixed and \(Q\le T^{2d}\). Let \(M\ge2\) be the threshold of Theorem~\ref{thm:short} for \(2d\) and tolerance \(1/2\). If \(Q\ge M\), then
\[
 \frac T2\le \left|T-\sum_{i<T}e(b^iw)\right|
      \le 2\pi\sum_{i<T}\circabs{b^iw}.
\]
For \(1<Q<M\), every \(b^iw\) still has reduced denominator \(Q\), hence distance at least \(1/Q\) to the nearest integer. In either case the mean distance is at least \(\gamma=\min(1/M,1/(4\pi))\). Apply the circle triangle inequality through \(f_i\) at each position.
\end{proof}

\begin{corollary}[Overlapping models]\label{cor:chain}
\evidence{Lean-supported paper proof}
Fix \(b,c,d\) as in Theorem~\ref{thm:rigidity}. For each integer \(j\ge0\) let \(T_j\ge1\) be an integer, and let \(y_j,z_j\in\Q\) have reduced denominators supported on \(c\) and at most \(T_j^d\). If for real observations \(f^{(j)}_i\), \(0\le i<T_j\), the left side of \eqref{eq:rigid} for this pair tends to zero as \(j\to\infty\), then \(y_j-z_j\in\Z\) for all sufficiently large \(j\).

In particular, consider consecutive blocks with rational models \(y_j,y_{j+1}\), whose start indices differ by an integer \(r_j\ge0\) and whose overlap has integer length \(T_j\ge1\). Suppose both denominators are supported on \(c\) and at most \(T_j^d\). If the combined mean error of the aligned models \(b^{r_j}y_j,y_{j+1}\) on their common observations tends to zero, then \(\den y_j\) is eventually constant.
\end{corollary}
\begin{proof}
Once the error is below \(\gamma\), \eqref{eq:rigid} forces integral difference. Multiplication by \(b^{r_j}\) preserves any denominator supported on \(c\), and an integer translation preserves it as well. Therefore \(\den y_{j+1}=\den y_j\) for all sufficiently large \(j\).
\end{proof}

Both the error and the size bound in this argument refer to the \emph{overlap}. Small mean error on a longer block can concentrate on its short shared part. A size bound polynomial in the original length suffices if that length is bounded by a fixed power of \(T\). Three examples distinguish the restrictions.

If a denominator shares a radix factor, the models \(0,1/b\) coincide after one shift. Their mean distance is \(1/(bT)\). Without the size bound, the fixed-support models \(0,c^{-T^2}\) have mean distance at most \(b^Tc^{-T^2}\to0\). Finally, for \(m\ge2\) put \(Q_m=b^m-1\) and take \(0,1/Q_m\) with \(T=mQ_m\). These denominators are prime to \(b\) and at most \(T\). Each length-\(m\) period has sum of the distances to zero at most \(\sum_{j<m} b^j/Q_m=1/(b-1)\). Their mean distance tends to zero, so arbitrary changes of prime support defeat uniform separation even with \(d=1\). The same example has first character mean tending to one. No fixed finite set of primes supports all these denominators since \(v_{Q_m}\le m\) and any fixed support would give \(Q_m\le C m\) in the proof of Theorem~\ref{thm:short}.

\section{A counting obstruction}\label{sec:count}
We bound the measure of points approximated on a prescribed block by any rational orbit whose denominator is bounded by a power of the block length. Let \(\lambda\) be Lebesgue probability measure on \([0,1)\).

\begin{theorem}[Polynomially many rational models]\label{thm:count}
\evidence{Lean-supported paper proof}
Fix integers \(b\ge2\) and \(d\ge0\), put \(a=b^2\), and set \(\eta_b=1/(16a^2)=1/(16b^4)\). For a prescribed integer start \(n\ge0\) and integer length \(T\ge2\), let \(E_{n,T}\) consist of all \(x\in[0,1)\) for which a rational \(y\) with \(\den y\le T^d\) satisfies \(T^{-1}\sum_{i<T}\circabs{b^{n+i}x-b^iy}\le\eta_b\). Then
\begin{equation}\label{eq:prob}
 \lambda(E_{n,T})\le T^{2d}\exp\!\left(-\frac{\lfloor T/2\rfloor}{4a}\right).
\end{equation}
Let \(T_k=B^k\) for integers \(k\ge1\) and a fixed integer \(B\ge2\), with at most \(O(T_k^A)\) prescribed integer starts at each scale \(k\ge1\), for a fixed real \(A\ge0\). For almost every \(x\in[0,1)\) there is \(k_0(x)\ge1\) such that none of the prescribed starts admits such a model at any scale \(k\ge k_0(x)\).
\end{theorem}
\begin{proof}
Fix \(y\). Under \(\lambda\), \(z=\{b^n x-y\}\) is uniform. Its first \(J=\lfloor T/2\rfloor\) base-\(a\) digits form an equiprobable word, since each cylinder has length \(a^{-J}\). An occurrence of digit \(1\) in position \(j+1\), \(0\le j<J\), puts \(\{b^{2j}z\}\) in \([1/a,2/a)\), at distance at least \(1/a\) from an integer since \(a\ge4\). The mean error condition allows at most \(K=\lfloor a\eta_b T\rfloor\) such occurrences.

Give digit \(1\) weight one and every other digit weight two. The total product weight of all words is \((2a-1)^J\), while each word with at most \(K\) occurrences has weight at least \(2^{J-K}\). Its class has probability at most
\[
 2^K(1-1/(2a))^J \le \exp(-J/(4a)).
\]
Indeed \(T\le3J\) for \(T\ge2\), hence \(K\log2\le 3J\log2/(16a)\le J/(4a)\). There are at most \(\sum_{q\le T^d}q\le T^{2d}\) rational models modulo one. A union bound gives \eqref{eq:prob}. Multiplying by the number of starts leaves a summable sequence in \(k\ge1\). The first Borel--Cantelli lemma gives the assertion without independence between starts.
\end{proof}
Borel--Cantelli also applies simultaneously to a countable family of prescribed schedules. In particular the integer parameters \(b,B,L,d\) of Theorem~\ref{thm:mean} give only countably many schedules, with \((LB-1)L^k\) starts per scale. On a late block the real tail \(b^jx-P_j=o(1)\) changes its mean error by \(o(1)\) uniformly for \(j\) beyond the block start. An arbitrarily accurate certificate for \(P\) would therefore give an approximation in \eqref{eq:prob}. For any fixed radix, almost every real is normal: the strong law on each of the finitely many residue classes of starting positions modulo a fixed word length gives that word's limiting frequency, and there are countably many words. Thus almost every normal real lacks these certificates. The measure estimate concerns schedules prescribed independently of \(x\).

\section{Countermodels for restricted data}\label{sec:pairs}

Digits of a binary stream \(s\) have zero-based indices, and
\[
  V(s)=\sum_{n\ge0}s(n)2^{-(n+1)}.
\]
When specifying positions of ones, we use the corresponding one-based indices. Let \(C\) be the binary Champernowne constant, given by the binary concatenation of the positive integers (the base-\(2\) analogue of Champernowne's decimal construction \cite{champernowne}), and \(F=\sum_{j\ge1}2^{-j!}\). Both have computable digits: find a digit of \(C\) by subtracting the lengths \(r2^{r-1}\) of complete groups of \(r\)-digit integers, then divide the offset by \(r\); for \(F\) test factorial membership. These searches can be bounded by the digit position.

For completeness, fix a binary word of length \(k\ge1\) and consider a prefix of \(C\) of length \(N\ge1\) reaching into the group of \(r\)-bit integers (\(r\ge1\)). For each integer \(1\le l\le r\), the prefix contains the first \(m_l\) complete \(l\)-bit integers, with \(m_l=2^{l-1}\) for \(l<r\) and \(0\le m_r\le 2^{r-1}\). When \(l\ge k\), an internal word with \(i\) bits to its right (\(0\le i\le l-k\)) has indicator periodic on all nonnegative integers with period \(2^{i+k}\) and mean \(2^{-k}\). Splitting the \(m_l\) integers into complete periods and a remainder of fewer than \(2^{i+k}\) integers gives count error at most \(2^{i+k}\) (including \(i=l-k\), where \(m_l\le 2^{l-1}<2^l=2^{i+k}\)). Summing over \(0\le i\le l-k\) bounds the internal count error of group \(l\) by \(2^{l+1}\). The remaining \((k-1)m_l\) starts in group \(l\), including words crossing a numeral boundary, contribute error at most \((k-1)m_l\le (k-1)2^{l-1}\). When \(l<k\), all \(l m_l\) starts in group \(l\) contribute at most \(l m_l\le(k-1)2^{l-1}\), using the boundary bound alone. Summing over \(1\le l\le r\) gives at most \(4\cdot 2^r+(k-1)2^r\), and the at most \(r\le 2^r\) starts in the last partial \(r\)-bit numeral add at most \(r\), yielding the explicit uniform error bound
\begin{equation}\label{eq:champ-err}
 \bigl|2^k A_C(N)-N\bigr|\le 2^k(k+5)2^r,
\end{equation}
where \(A_C(N)\) is the number of occurrences of the word among the first \(N\) starts. A prefix reaching group \(r\ge2\) contains all \((r-1)2^{r-2}\) bits of group \(r-1\), so \(N\ge(r-1)2^{r-2}\). Dividing \eqref{eq:champ-err} by \(2^k N\) bounds the normalized error by \(4(k+5)/(r-1)\to0\). Thus \(C\) is normal. For \(F\), at most \(1+\lfloor\log_2 N\rfloor\) digits through \(N\) are ones. Each one meets at most \(k\) windows of length \(k\). Consequently \(0^k\) has ordinary frequency one for every fixed \(k\), which implies \(D_N^{(2)}(F)\to1\) on testing \([0,2^{-k})\) and letting \(k\) grow. The same reasoning works for any canonical stream with zero density of ones.

\begin{proposition}[Shared odd-primary data]\label{prop:odd}
\evidence{Paper proof}
There are computable reals with discrepancy limits zero and one and specified computable rational models for their shifts that agree in every odd-primary component at every cutoff, have identical reduced odd denominators, and each approximate its shift to absolute error less than \(2^{-n}\).
\end{proposition}
\begin{proof}
Use \(x=C,F\), write \(t_n(x)=\{2^nx\}\), and let \(R_n=2^n\sum_{k=1}^n 2^{-k}/k\) with \(R_0=0\). Set
\[
 Y_n(x)=R_n+2^{-n}\big(\lfloor2^nt_n(x)\rfloor-\lfloor2^nR_n\rfloor\big).
\]
The two values \(Y_n(x)\) and \(t_n(x)\) lie in the same half-open mesh cell of width \(2^{-n}\), and \(\lfloor2^nt_n(x)\rfloor\) uses only \(n\) digits. The difference \(Y_n(x)-R_n\) is dyadic. In \(\Q/\Z\) the primary components are the summands in the decomposition into elements of prime-power order, obtained by the Chinese remainder theorem on any common denominator. A dyadic difference changes only the component of order a power of two. The order of the odd part is the odd part of the reduced denominator, giving both equalities.
\end{proof}

\begin{proposition}[Logarithmic word laws]\label{prop:log}
\evidence{Paper proof}
There are computable binary expansions agreeing in the limiting logarithmic frequency \(2^{-k}\) of every word of length \(k\ge1\). Both contain every word infinitely often. One has ordinary discrepancy tending to zero and the other has ordinary upper limiting discrepancy one.
\end{proposition}
\begin{proof}
Keep \(C\) as the first expansion. For the second, replace its digits at the one-based positions in \([M_j,jM_j)\) by zero, where \(M_j=2^{2^j}\), \(j\ge2\). Keep all other digits in place. These are computable disjoint intervals. The logarithmic frequency weights the start at position \(n\) by \(1/n\) and divides the weighted count up to \(N\) by \(H_N=\sum_{n=1}^N1/n\). Partial summation turns \(A(n)=pn+o(n)\) for an ordinary count into
\[
 A(N)/N+\sum_{n<N} A(n)/(n(n+1))=pH_N+o(H_N).
\]
An edited start for a \(k\)-word must belong to some \([M_j-k+1,jM_j)\). For all late \(j\) this interval has harmonic mass at most \(\log j+O_k(1)\). If \(J\) is the last index whose enlarged interval has begun by \(N\), the total charge is \(O_k(J\log J)\) and \(\log N\gg 2^J\). The normalized charge tends to zero, including at cutoffs inside an edited interval. Positive limiting logarithmic frequencies prove infinite occurrence and ensure the expansion is canonical. At \(N=jM_j\), at least \((j-1)M_j-k+1\) starts see \(0^k\), so the upper limiting discrepancy is at least \(1-2^{-k}\) for every \(k\), and hence is one.
\end{proof}

Any finite collection of observations consisting of counts of bounded-length words through specified finite cutoffs depends on a single finite prefix. Attach either tail \(C\) or \(F\) to that prefix to preserve every observation while obtaining discrepancy limits zero and one: after the prefix the two orbits are exactly those of the two tails with a fixed delay. Even knowing the exact limiting frequencies through a fixed word length \(k\) is insufficient. A periodically repeated binary de Bruijn cycle of order \(k\) has the same frequencies as \(C\) at every length at most \(k\), and some longer word is absent. Such a cycle is obtained by an Euler tour of the directed graph of \((k-1)\)-words, with one edge for each \(k\)-word joining its prefix to its suffix. The graph is strongly connected by appending the target word, and every vertex has indegree and outdegree two. The order-one cycle is \(01\).

\section{Eight limits of particular methods}\label{sec:methods}

Write \(D^+(x)=\limsup_N D_N^{(2)}(x)\). The examples below isolate the information lost by each method and give a comparison or a sufficient replacement. The arguments apply to the stated statistics and model classes.

\begin{enumerate}[leftmargin=*,label=\textbf{\arabic*.}]
\item \textbf{Fixed-support certificates.}
The binary Stoneham number \(x=\sum_{j\ge1}(3^j2^{3^j})^{-1}\) (originally constructed by Stoneham \cite{stoneham}, with normality also obtained via Bailey and Crandall \cite[Theorem~4.8, p.~535]{bailey}) illustrates \cref{thm:mean}. For its shifted prefix \(P_n\) retaining terms with \(3^j\le n\), the last insertion gives reduced denominator \(3^j\), since its numerator at that denominator is a unit modulo \(3\); the earlier terms have smaller denominators. The omitted shifted tail is bounded by \(1/(3^{j+1})\), since the remaining exponents relative to \(n\) are distinct positive integers. Between insertions the models are exact radix orbits. On the schedule \(B=L=2\), a shell has at most two insertions, affecting at most two block interiors. On the other blocks \(n/3<\den P_n\le n\), and \(n<4\cdot4^k\le (2^k)^3\) at sufficiently large \(k\). In the full Taylor expansion of \(\ln 2\), the shifted truncations instead have odd denominators whose logarithm is \(n+o(n)\) in the sense of item~3. On this geometric schedule \(\log T=O(\log n)\), so these particular truncations violate every fixed polynomial cap. A different family of accurately shadowing models, or estimates uniform as the prime support grows, would be needed to use this approach.

\item \textbf{Overlaps and prescribed starts.}
By \cref{cor:chain}, vanishing error on actual overlaps with fixed support and the polynomial cap forces a constant denominator along an eventual chain. By \cref{thm:count}, even among normal reals the same fixed schedule typically has no models meeting the mean-error threshold \(1/(16b^4)\) at all on late blocks. The disjoint Stoneham blocks from item~1 give the comparison: their starts are fixed on the schedule and their denominators grow without requiring agreement across insertion boundaries. When using a different schedule, \cref{prop:cover} identifies the required coverage and within-block accuracy.

\item \textbf{Nearby Taylor cutoffs.}
Put \(T_j=\sum_{\ell=1}^j2^{-\ell}/\ell\). Let \(h=o(K)\) be a nonnegative integer, and form any integer affine combination of \(T_j\) with \(K\le j\le K+h\). Its reduced odd denominator \(q\) satisfies \(\log q=K+o(K)\) uniformly in the combination, also after any binary shift. To see this, fix \(m\ge1\). For primes \((K+h)/(m+1)<p\le K/m\), each truncation has exactly the multiples \(p,\ldots,mp\) below its cutoff, and \(p^2>K+h\) for large \(K\). The residue of \(pT_j\) modulo \(p\) is \(\sum_{\ell=1}^m2^{-\ell}/\ell\), by Fermat's little theorem. This fixed positive rational is a unit for all sufficiently large \(p\) and the residue survives the affine combination. By the prime number theorem in the form \(\sum_{p\le y}\log p\sim y\) \cite{hadamard}, any fixed number \(J\) of these disjoint bands gives \(\liminf(\log q)/K\ge 1-1/(J+1)\). The upper bound is \(\log\operatorname{lcm}(1,\ldots,K+h)=K+o(K)\), since the prime powers beyond the first powers contribute \(O(\sqrt{K+h}\log^2(K+h))\). The residue argument is independent of the coefficients and the count of truncations. At width \(K\), the odd-primary projections of the increments \(2^{-j}/j=T_j-T_{j-1}\) for \(K<j\le2K\) generate the cyclic group with order the odd part of \(\operatorname{lcm}(1,\ldots,2K)\). Each maximal prime power divides a multiple in this range, and the order of a generated subgroup is the least common multiple of the generator orders. Thus adding integer multiples of these increments to \(T_K\) can cancel its odd part, producing an integer affine combination with odd denominator one. This construction supplies no bound on the coefficients, the dyadic denominator or the propagated error.

\item \textbf{Balanced trace sources.}
Higher degree can describe exactly the same finite cutoff. For example, take all roots \(\alpha\) of \(X^r-2^a\), with integers \(a\ge0,\ r\ge1\), and an integer \(t>a/r\). At every real cutoff \(H>0\),
\[
 \sum_{(t-a/r)j\le H,\ j\ge1}\frac{\sum_\alpha\alpha^j}{j\,2^{tj}}
 =\sum_{(rt-a)m\le H,\ m\ge1}\frac{2^{-(rt-a)m}}m.
\]
The trace vanishes unless \(r\mid j\); at \(j=rm\) it is \(r2^{am}\), which gives the identity including the cutoff. In particular, the degree-\(r\) polynomial \(X^r-2^{r-1}\) at \(t=1\) recovers the original Taylor truncations for \(\ln 2\) with their full odd-prime obstruction. As an equal-modulus comparison, \(X^2-1\) and \(X^2+1\) at \(t=1\) both have roots of modulus one, yet give respectively the nonalternating sum \(\sum_{2m\le H,\,m\ge1}2^{-2m}/m\) and the alternating sum \(\sum_{2m\le H,\,m\ge1}(-1)^m 2^{-2m}/m\), by summing the powers of \(\pm1\) and \(\pm i\). A change in the actual finite rational truncations, together with a quantitative precision bound, is required to improve a certificate. This calculation treats only these explicit sources.

\item \textbf{Prime-by-prime statistics.}
\Cref{prop:odd} preserves the entire odd-primary vector at every cutoff while permitting opposite discrepancy limits. Per-position shadowing alone allows the dyadic corrections in its proof. It gives no control over successive displacements of the rational models or the characters of the full phase. For a positive criterion, cancellation for every fixed nonzero frequency of the full shadowing model implies normality when the mean circle error tends to zero, by the Lipschitz bound and Weyl's criterion.

\item \textbf{Higher-rank density and fixed projections.}
Let \(E(t)=1-\cos(2\pi t)\) and prescribe positive integers \(M_N\) with \(\log M_N=o(N)\). There is an irrational \(x\) with \(D_N^{(2)}(x)\to1\) and
\[
  \frac1N\sum_{n<N}\max_{1\le m\le M_N} E(m2^nx)\ \longrightarrow\ 0 .
\]
Put \(L_N=\lceil\log_2 M_N\rceil+\lceil\log_2(N+1)\rceil\) and choose one-positions \(p_j\) with \(p_{j+1}\ge 2p_j\), sufficiently late that \(L_N/N\le j^{-2}\) whenever \(N\ge p_j/2\). For large \(N\) with \(N+L_N\le 2N\), let \(j\) count the ones through \(2N\). Then at most \(jL_N\le N/j\) starts may see a one in their next \(L_N\) places. At other starts every allowed multiple is at circle distance at most \(1/(N+1)\), implying \(E\le2\pi^2/(N+1)^2\). The one density is zero, so the discrepancy limit is one. The increasing gaps also ensure irrationality. Furstenberg's topological theorem gives density of \(\{2^u3^v x\bmod1:u,v\ge0\}\) for this \(x\) and for the normal irrational \(C\) \cite[Theorem~IV.1, p.~48]{furstenberg}. A positive lower mean at bounded multiplier cost would force a positive lower binary mean by \(E(mt)\le m^2 E(t)\), which follows by factoring \(1-e(mt)\) and taking squared norms. The fixed-projection issue is separate: individually normal approximants may have a nonnormal limit, even for \(x_j=2^{-j}C\to0\). To pass from fixed-width projections to one full orbit requires estimates uniform in the cutoff and an error bound as width grows.

\item \textbf{Complex products at the endpoint.}
Write \(P_N(a,z)=\prod_{n<N}(1+az^{2^n})\) and \(z=e(x)\). For every prescribed \(0<a_N<1\) tending to one, an irrational \(x\) has discrepancy tending to one and both \(N^{-1}\log|P_N(a_N,z)|\) and \(N^{-1}\log|P_N(1,z)|\) tending to zero. Set \(d_N=1-a_N\) and
\[
 h(k)=\max\big(2^{-\sqrt{k}},\sqrt{\sup_{m\ge k}d_m}\big).
\]
Starting sufficiently late with \(h(p_1)\le1/4\), take ones at \(p_{j+1}=p_j+\lfloor\log_2(1/h(p_j))\rfloor\). For \(p_j\le k<p_{j+1}\) the shifted fractional part lies between \(h(k)\) and \(2/3\): the next one is at most \(\lfloor\log_2(1/h(p_j))\rfloor\) places away and later gaps are at least two. Thus \(\circabs{2^k x}\ge h(k)\). The gaps tend to infinity, so \(x\) is irrational with one density zero. Telescoping gives \((1-z)P_N(1,z)=1-z^{2^N}\), whose numerator has magnitude between \(4\cdot2^{-\sqrt N}\) and two for large \(N\). For \(u=e(2^n x)\) with \(p_1\le n+1\le N\), factoring \(1-u^2\) gives \(|1+u|\ge2h(n+1)\ge2\sqrt{d_N}\). Eventually that lower bound holds for the finitely many earlier factors as well. The squared modulus ratio \(|1+a_Nu|^2/|1+u|^2=a_N+d_N^2/|1+u|^2\) tends to one uniformly for \(n<N\). This proves the moving limit.

 For comparison, at any fixed \(0<a<1\), the per-factor deficit \(\log(1+a)-\log|1+ae(t)|=-\frac12\log(1-2a E(t)/(1+a)^2)\) is bounded above and below by positive multiples of \(E(t)\). Hence its mean tends to zero for this sparse example. A strictly positive lower mean of this fixed-parameter deficit is equivalent to \(D^+(x)<1\). Indeed, a subsequence with mean \(E\to0\) concentrates near zero on the circle, and testing the complement of a small arc there gives upper limiting discrepancy one. Conversely, along a subsequence with discrepancy tending to one, an interval or its circular complement of length tending to zero contains a proportion tending to one. If its center is \(u\) and its circle radius \(r\to0\), some two consecutive shifts both lie in that arc, by counting bad starts. Doubling then gives \(\circabs{u}\le3r\). The empirical measures concentrate near zero along the subsequence, and the mean of \(E\) tends to zero.

\item \textbf{Irrationality exponents and digit runs.}
If \(x\) has irrationality exponent two, a run from position \(m+1\) through \(r\) gives error at most \(2^{-r}\) from a dyadic of denominator at most \(2^m\) (round up for ones). For any \(\epsilon>0\), the exponent assumption forbids \(r>(2+\epsilon)m\) at all sufficiently late starts: the approximating reduced denominators tend to infinity as \(m\to\infty\). Iterating the inequality at successive runs gives at least order \(\log N\) transitions through \(N\). This order is attained by \(K=\sum_{j\ge0}2^{-2^j}\); its ones have zero density and there are only \(O(\log N)\) transitions. The continued fraction of \(2K\) has bounded partial quotients \cite[pp.~240--242]{folded}, so \(K\) has exponent two as well. Even occurrence of every word is compatible with zero one density: insert an enumeration of all binary words, ordered by length, at positions \(2^{2^j}\) in an otherwise zero stream. The total length of the first \(j\) words is \(O(j\log(j+1))\), giving zero density including within the last inserted word. In contrast, positive lower ordinary transition density forces a positive lower mean of \(E\), because a transition corresponds to a shift in \([1/4,3/4)\). Item~7 then gives \(D^+<1\). Normality requires the ordinary limiting frequencies at every resolution.
\end{enumerate}

\section{Computability on digit-program names}\label{sec:computable}

For a partial-recursive program code \(p\), write \(\varphi_p(n)\) for its partial value on \(n\). In Lean the type of codes is \nolinkurl{Nat.Partrec.Code}. A code \emph{names} \(s:\N\to\{0,1\}\) if \(\varphi_p(n)\) is defined and equal to \(s(n)\) for every integer \(n\ge0\). In \lean{ComputableNormalityUndecidability}{393}{\nolinkurl{Names}}, the stream is Bool-valued and the exact condition is
\(\texttt{code.eval n = Part.some (s n).toNat}\) for all \(n\).
The promise \lean{ComputableNormalityUndecidability}{398}{\nolinkurl{TranscendentalZeroOnePromise}} on a stream \(s\) is
\[
 V(s)\text{ is transcendental over }\Q,\qquad
 \big(D_N^{(2)}(V(s))\to0\big)\ \text{or}\ \big(D_N^{(2)}(V(s))\to1\big).
\]
In particular these are limits of the full sequence of discrepancies. The predicate contains no modulus of convergence. Transcendence ensures \(0<V(s)<1\) and removes binary expansion ambiguity. A recognizer is a computably enumerable (c.e.) predicate on codes, required to agree with the desired answer on every promised name. Its behavior outside that domain is unrestricted. A decider in the formal theorem is a computable predicate on codes with the same specification.

\begin{lemma}[Finite edits]\label{lem:edit}
\evidence{Lean-supported paper proof}
If \(s(n)=t(n)\) for all \(n\ge m\), then \(V(s)-V(t)=j/2^m\) for some integer \(j\). Adding such a dyadic rational to any real preserves transcendence, binary normality, and convergence of binary discrepancy to any given value, in both directions.
\end{lemma}
\begin{proof}
The tails cancel in the two absolutely convergent series. Multiplication of their finite difference by \(2^m\) gives an integer. For the translated real, every shift at index at least \(m\) is identical, so the discrepancy difference is bounded by \(m/N\) for \(N>0\), uniformly over test intervals. Translation by a rational preserves algebraicity in both directions by substituting into a polynomial, applying the one-way rational translation statement to both \(+q\) and \(-q\). Each of the four formal components is individually kernel-verified in
\lean{ComputableNormalityUndecidability}{57}{\nolinkurl{binarySeries_finite_edit_dyadic}},
\lean{ComputableNormalityUndecidability}{85}{\nolinkurl{transcendental_add_rat}} (one-way),
\lean{ComputableNormalityUndecidability}{177}{\nolinkurl{binaryDiscrepancy_tendsto_iff_add_dyadic}}
and \lean{ComputableNormalityUndecidability}{211}{\nolinkurl{normalIn_two_iff_add_dyadic}} (via \lean{ComputableNormalityUndecidability}{101}{\nolinkurl{shift_add_dyadic_eq}} and \lean{ComputableNormalityUndecidability}{155}{\nolinkurl{binaryDiscrepancy_sub_le_of_eventually_eq}}); the two-way transcendence consequence and bundling into one lemma are written deductions.
\end{proof}

\begin{proposition}[Two computable transcendental endpoints]\label{prop:endpoints}
\evidence{Paper proof}
There are computable binary streams \(a,b_*\) whose values are transcendental, with \(V(a)\) binary normal and \(D_N^{(2)}(V(b_*))\to1\).
\end{proposition}
\begin{proof}
For \(b_*\) use the digits of \(F\) from \cref{sec:pairs}. The expansion is irrational because it has unbounded gaps between ones and infinitely many ones. The truncation at digit \(j!\) has reduced denominator \(q=2^{j!}\) and positive error at most \(2q^{-(j+1)}\). The elementary transcendence test used here is as follows. For an algebraic irrational \(\alpha\) with irreducible integer polynomial of degree \(d\), the polynomial at a rational \(p/q\) has magnitude at least \(q^{-d}\), since it is nonzero. A derivative bound in a neighborhood of \(\alpha\) gives \(|\alpha-p/q|\gg_\alpha q^{-d}\). The factorial truncations violate this for every fixed \(d\).

For \(a\), let \(U_j\) concatenate all \(2^j\) binary words of length \(j\) in lexicographic order, with \(\ell_j=j2^j\). Starting at \(S_0=0\), append at stage \(j\ge1\) the smallest positive number of complete copies of \(U_j\) that gives total length \(S_j\ge(j+1)^2(S_{j-1}+\ell_j)\). The digit is computable by finding its stage and reducing its offset modulo \(\ell_j\). For a \(k\)-word with \(k\le j\), a complete copy has \((j-k+1)2^{j-k}\) occurrences internal to the length-\(j\) words, with \((k-1)2^j\) other starts, including the boundary of the copy. Its error from \(\ell_j2^{-k}\) is therefore at most \((k-1)2^j\), independently of the continuation. Completed late copies have uniformly small relative error for fixed \(k\). At a cutoff in stage \(j\ge2\), the unfinished copy contributes at most \(\ell_j\), and \(\ell_j/S_{j-1}\le 2/[j(j-1)]\). The finite initial stages have vanishing weight. This proves normality at every cutoff, hence irrationality. Compare \(V(a)\) to the rational with the first \(S_{j-1}\) digits unchanged followed by \(U_j\) repeated periodically. The rational has reduced denominator \(q_j\le2^{S_{j-1}}(2^{\ell_j}-1)\) and agrees for \(S_j\) digits. Its error is nonzero and at most \(2^{-S_j}\le q_j^{-(j+1)^2}\). The denominators tend to infinity because these rationals converge to an irrational. The same polynomial test gives transcendence.
\end{proof}

\begin{theorem}[Two uniform halting reductions]\label{thm:switch}
\evidence{Lean-supported paper proof}
Fix streams \(a,b_*\) with the transcendence, normality and discrepancy properties in \cref{prop:endpoints}. For a code \(c\) and index \(n\), set \(S_{a,b_*}(c,n)=b_*(n)\) if its bounded evaluation at bound \(n\) on input \(0\) succeeds, and set it to \(a(n)\) otherwise. The stream \(S_{b_*,a}(c,\cdot)\) reverses the endpoints. Both streams satisfy the promise for every \(c\). The first is normal exactly when \(\varphi_c(0)\) is undefined, and the second exactly when it is defined. If \(a,b_*\) are computable, both switches are jointly computable in \((c,n)\) and there are total computable maps from \(c\) to codes naming the corresponding streams.
\end{theorem}
\begin{proof}
Use the bounded evaluator \texttt{Nat.Partrec.Code.evaln n c 0}, which is computable, sound, complete as the bound grows, and monotone in successful evaluations. If \(\varphi_c(0)\) is undefined the switch keeps its first stream. If it is defined there is a bound \(m\) from which the switch always uses its second stream at the original indices. Apply \cref{lem:edit} in each case. Bounded evaluation and the two computable endpoints give joint computability. A code for the two-argument output, on an encoded pair, can be specialized to \(c\) by the computable currying operation on codes (the \(s\)-\(m\)-\(n\) construction).

The four individually kernel-verified declarations under the endpoint hypotheses are \lean{ComputableNormalityUndecidability}{239}{\nolinkurl{switched_real_promise_and_reduction}},
\lean{ComputableNormalityUndecidability}{295}{\nolinkurl{reversed_switched_real_reduction}},
\lean{ComputableNormalityUndecidability}{223}{\nolinkurl{switchedDigit_computable}}, and
\lean{ComputableNormalityUndecidability}{406}{\nolinkurl{exists_switchedDigit_index_map}}; applying the latter two to both \((a,b_*)\) and \((b_*,a)\) and bundling the clauses into one theorem statement are written deductions.
\end{proof}

\begin{theorem}[Nonrecognizability on names]\label{thm:names}
\evidence{Lean-supported paper proof}
Assume \(a,b_*\) are computable streams with the properties in \cref{prop:endpoints}. There is no computable predicate on codes agreeing with binary normality on all promised names. There is no c.e.\ predicate on codes agreeing with binary normality on all promised names, and none agreeing with binary nonnormality on all promised names.
\end{theorem}
\begin{proof}
Let \(K_0=\{c:\varphi_c(0)\downarrow\}\). The halting set is c.e.; its complement is not c.e., since dovetailing two recognizers would decide \(K_0\). This would decide diagonal halting by specializing a code to its own input, contrary to the diagonal program that halts precisely when the decider predicts divergence. Pull a hypothetical normality recognizer back along the first computable name map of \cref{thm:switch}. The result recognizes the complement of \(K_0\), since every output of the map is a promised name. For nonnormality pull back along the reverse map to get the same contradiction. A computable predicate is c.e.

The three exclusions are individually kernel-verified under the endpoint hypotheses in \lean{ComputableNormalityUndecidability}{468}{\nolinkurl{no_computable_normality_decider_on_names}},
\lean{ComputableNormalityUndecidability}{428}{\nolinkurl{no_re_normality_recognizer_on_names}}, and \lean{ComputableNormalityUndecidability}{448}{\nolinkurl{no_re_nonnormality_recognizer_on_names}} (using the direct switched-family versions \lean{ComputableNormalityUndecidability}{351}{\nolinkurl{normality_not_computable}}, \lean{ComputableNormalityUndecidability}{367}{\nolinkurl{normality_not_re}}, and \lean{ComputableNormalityUndecidability}{380}{\nolinkurl{normality_not_co_re}}); bundling the three separately verified declarations into one theorem is a written step.
In each declaration, a predicate \(P\) is quantified over \emph{all} codes, with its correctness clause quantified over all pairs \((p,s)\) such that \(p\) names \(s\) and \(s\) satisfies the promise.
\end{proof}

\begin{corollary}[Explicit promise problem]\label{cor:instantiate}
\evidence{Lean-supported paper proof}
Binary normality on promised names is neither decidable, c.e., nor co-c.e. Even a partial algorithm required to return a correct yes/no answer and halt on every promised name is impossible.
\end{corollary}
\begin{proof}
Instantiate the three individually kernel-verified exclusions of \cref{thm:names} by the written endpoint existence proof of \cref{prop:endpoints}. The partial-decider consequence is a written deduction: if a partial computable decider \(D\) halts on every promised name and outputs \(1\) exactly on normal promised names, then the predicate \(p\mapsto(D(p)\downarrow 1)\) is a c.e.\ predicate on codes agreeing with normality on the promise, contradicting \cref{thm:names}.
\end{proof}

Ki and Linton prove boldface Borel \(\boldsymbol{\Pi}^0_3\)-completeness for fixed-base normality \cite{kilinton}. Becher, Heiber and Slaman construct an effective reduction to normality in all bases versus in no base, establishing both effective \(\Pi^0_3\)-completeness and boldface Borel \(\boldsymbol{\Pi}^0_3\)-completeness for absolute normality and for each fixed base \cite[Theorem~3.2 and Corollaries~3.3--3.4, p.~67; proofs p.~76]{bhs}; their reduction separates absolute normality from absolute abnormality without prescribing a discrepancy limit. Becher and Slaman prove index (\(\Sigma^0_4\)) and boldface Borel (\(\boldsymbol{\Sigma}^0_4\)) completeness for normality to at least one base \cite[Theorem~1]{becherslaman}. Those theorems classify unrestricted sets of reals or computable-real indices in the Borel and arithmetical hierarchies rather than recognition on the narrower transcendental \(\{0,1\}\)-discrepancy promise treated here. The present construction uses a single bounded switch in each direction to enforce the two discrepancy limits and transcendence on every output; it gives no completeness classification.

Rice's theorem concerns nontrivial index properties of c.e.\ sets \cite{rice}. The quantification here is over total bit names with an additional analytic promise. In particular, the reductions enforce this promise even on nonhalting inputs. Specker's computable increasing rational sequence with no computable convergence modulus \cite{specker} provides precedent for separating computable finite data from effective limiting information. Here each value is a computable real, while the obstruction concerns the limit of an orbit statistic.

Constructing a particular normal number has different quantifiers. Sierpi\'nski's 1917 effective determination \cite{sierpinski}, its computable reformulation \cite{becherfigueira}, and Turing's unpublished construction as transcribed and corrected by Becher, Figueira and Picchi \cite{turing} address existence and generation. Becher, Heiber and Slaman construct a computable absolutely normal Liouville number \cite[p.~2939]{bhsliouville}. Lutz and Mayordomo give an online prefix-generation algorithm, first posted in 2016, that outputs the first \(n\) binary digits of an absolutely normal real in time \(n\,\mathrm{polylog}(n)\) \cite[Theorem~5.5, arXiv p.~17]{lutz}.

\section{Independence for theory-dependent descriptions}\label{sec:rosser}

For a primitive recursive digit program \(s\), let \(\mathsf{Norm}(s)\) be the arithmetical sentence that for every coded word of length \(k\ge1\) and every \(m\ge1\) there is \(N_0\ge1\) such that, for \(N\ge N_0\), the number \(A_s(N)\) of occurrences of that word among the first \(N\) starts satisfies
\[
 m\,|2^k A_s(N)-N| < 2^k N .
\]
In PA, exponentiation \(2^x\), factorial \(j!\), the integer logarithm \(\lfloor\log_2 N\rfloor\), the digit functions of \(C\) and \(F\), the bounded-search switch \(s_T\) below, and the word-count function \(A_s(N)\) are total primitive recursive functions represented by \(\Sigma_1\) defining formulas (via G\"odel's \(\beta\)-function) whose totality and defining recursions are provable in PA. Use the usual interpretation of arithmetic in PA and in ZFC.

\begin{theorem}[Rosser switch]\label{thm:rosser}
\evidence{Paper proof}
For each \(T\) equal to PA or ZFC, there is an effectively specified primitive recursive stream \(s_T\) such that if \(T\) is consistent, neither \(\mathsf{Norm}(s_T)\) nor its negation is provable in \(T\). In the standard natural numbers under that consistency hypothesis, \(s_T\) is the unmodified binary concatenation stream of \(C\).
\end{theorem}
\begin{proof}
Use a primitive recursive proof test \(\mathsf{Prf}_T(p,\sigma)\) with enumeration stages included in the code if necessary. Use diagonal substitution to construct a sentence \(R\) with
\[
 T\vdash R\ \longleftrightarrow\ \forall p\,\neg W(p),
 \qquad W(p):=\mathsf{Prf}_T(p,\ulcorner R\urcorner)
       \ \land\ \forall q\le p\,\neg\mathsf{Prf}_T(q,\ulcorner\neg R\urcorner).
\]
This is the Rosser fixed point \cite[Theorem~II]{rosser}. The digit at zero-based index \(n\) is the original digit of \(C\) if no \(p\le n\) passes \(W\), and the digit of \(F\) at the same index otherwise. This is primitive recursion, with a bounded search and a permanent switch.

The needed implication between proofs is internal to \(T\). The explicit integer-only count establishing \eqref{eq:champ-err} in \cref{sec:pairs} formalizes uniformly in PA for every coded word of length \(k\ge1\) and tolerance denominator \(m\ge1\): for each complete or partial group of \(l\)-bit integers (\(1\le l\le r\)), division with remainder bounds \(|2^k A_{l,i}-m_l|\le 2^k\cdot 2^{i+k}\) per internal offset \(0\le i\le l-k\), summing to \(2^k\cdot 2^{l+1}\) for \(l\ge k\), and bounds the remaining (or short-group, when \(l<k\)) starts by \(2^k(k-1)2^{l-1}\). Summing over \(1\le l\le r\) and adding at most \(2^k r\le 2^k 2^r\) for the final partial numeral yields \(|2^k A_C(N)-N|\le 2^k(k+5)2^r\) by finite induction in PA. Since any prefix reaching group \(r\ge2\) satisfies \(N\ge(r-1)2^{r-2}\), choosing \(N_0\) beyond the end of group \(4m(k+5)+1\) ensures \(r-1>4m(k+5)\) and hence \(m|2^k A_C(N)-N|\le 2^k\cdot 4m(k+5)2^{r-2}<2^k N\) for all \(N\ge N_0\). If no \(p\) passes \(W\), every digit of \(s_T\) agrees with \(C\), so PA proves \(\mathsf{Norm}(s_T)\). If \(W(p)\) holds, at most \(p\) initial digits differ from \(F\); induction gives \(j!\ge 2^{j-1}\) for \(j\ge1\), so for the one-bit word \(1\) (\(k=1\)) the count satisfies \(A_{s_T}(N)\le p+1+\lfloor\log_2 N\rfloor\). For every \(N\ge\max(64,8(p+2))\), letting \(r=\lfloor\log_2 N\rfloor\ge6\), induction on \(r\ge6\) gives \(8r\le 2^r\le N\) while \(8(p+1)<8(p+2)\le N\); hence \(8 A_{s_T}(N)\le 8(p+1+r)<2N\), so \(4 A_{s_T}(N)<N\) and \(4|2A_{s_T}(N)-N|>2N\), contradicting \(\mathsf{Norm}(s_T)\) with \(k=1,m=4\). Thus \(T\vdash\mathsf{Norm}(s_T)\leftrightarrow R\), using the same PA argument in the natural numbers of ZFC.

For consistency alone to suffice, compare proof codes. If there is an actual proof of \(R\) with code \(p\), consistency rules out any actual proof of \(\neg R\). The true finite checks establishing \(W(\bar p)\) are provable in \(T\) and imply \(\neg R\). If there is an actual proof of \(\neg R\) with code \(q\), consistency rules out proof codes of \(R\) below \(q\). The theory verifies this finite list, and for any purported proof code \(p\ge q\) the fixed code \(q\) defeats \(W(p)\). It follows in \(T\) that \(\forall p\,\neg W(p)\), implying \(R\). Either case contradicts consistency. Externally there is consequently no proof of \(R\), so the switch never fires.
\end{proof}

\begin{remark}[Scope, weak arithmetic, and classical constants]\label{rem:classical}
The independent sentence \(\mathsf{Norm}(s_T)\) refers to the specified theory-dependent digit program \(s_T\). In the standard model under consistency, \(s_T\) denotes the same real number as \(C\); a theory proving equality of those two digit descriptions would transfer the elementary normality proof for \(C\). The argument is consistency-relative, the frequency transfer is a written arithmetization argument, and no internal transcendence theorem is used. While interpreting Robinson arithmetic \(Q\) suffices for the syntactic Rosser step, transferring independence to \(\mathsf{Norm}(s_T)\) requires a proof in \(T\) of the equivalence \(\mathsf{Norm}(s_T)\leftrightarrow R\); in very weak bounded arithmetic (such as \(I\Delta_0+\neg\mathrm{Exp}_2\), where \(\mathrm{Exp}_2\) is the totality sentence \(\forall x\,\exists y\,(y=2^x)\), under prefix-code encodings requiring powers of two to exist), coding validity can fail independently of the Rosser fixed point.

The results in this section and elsewhere in the manuscript establish neither normality nor independence for any named classical constant, including \(\ln 2\), \(\pi\), \(e\), \(\sqrt 2\), and the Erd\H{o}s--Borwein constant; the normality of each of those constants remains open.
\end{remark}

\section{Formal verification and data availability}
The \repo{NormalityChecker}{Lean source} is pinned to \texttt{e4809cf} (registry entries \texttt{AC-01}, \texttt{AC-02}, \texttt{AC-03}, \texttt{AC-04}, \texttt{AC-10}, and \texttt{AC-11}). Two complete statements in this manuscript match a single pinned declaration at \texttt{e4809cf} (\emph{Kernel-verified}): \cref{thm:short} (\texttt{AC-04}) is \lean{SmoothOrbit}{207}{\nolinkurl{fixed_support_orbit_estimate}}, and \cref{thm:rigidity} (\texttt{AC-03}) is \lean{OrbitRigidity}{32}{\nolinkurl{smooth_overlap_separation}}. The remaining formal links support \emph{Lean-supported paper proof} statements. In \cref{thm:mean} (\texttt{AC-01}), \lean{ArithmeticCertificate}{62}{\nolinkurl{normal_of_approximate_arithmetic_certificate}} takes the short-orbit estimate as an input, which \lean{SmoothOrbit}{207}{\nolinkurl{fixed_support_orbit_estimate}} supplies directly for every coprime pair \((b,c)\). The centered-count estimates underlying \cref{prop:cover} (\texttt{AC-02}) and the fixed-resolution bound are formally specialized to radix observations in \lean{QualityCoverage}{13}{\nolinkurl{Coverage.fixed_quality_cover}}, \lean{QualityCoverage}{38}{\nolinkurl{Coverage.asymptotic_fixed_quality}}, \lean{PartialUniformity}{75}{\nolinkurl{discrepancy_of_uniform_share}} and \lean{PartialUniformity}{39}{\nolinkurl{discrepancy_of_fixed_cells}}. Both halves of \cref{cor:chain} (\texttt{AC-03}) are individually verified in \lean{OrbitRigidity}{83}{\nolinkurl{smooth_overlap_rigidity}} and \lean{OrbitRigidity}{101}{\nolinkurl{overlap_chain_denominator}}. In \cref{thm:count} (\texttt{AC-10}), Lean verifies only the finite weighted word inequality \lean{ModelCounting}{14}{\nolinkurl{few_symbol_words}}; uniformity of the shifted point, the cylinder probability calculation, the union bound over rational models, and the Borel--Cantelli deduction are written arguments. In \cref{sec:computable} (\texttt{AC-11}), each individual finite-edit, switching, parameter-specialization, and nonrecognizability component cited in \cref{lem:edit,thm:switch,thm:names} is kernel-verified in \repo{NormalityChecker/ComputableNormalityUndecidability.lean}{\nolinkurl{ComputableNormalityUndecidability.lean}} under explicit endpoint hypotheses, while the two-way transcendence deduction, bundling of components, endpoint existence (\cref{prop:endpoints}), and the partial-decider consequence (\cref{cor:instantiate}) are written arguments, as are the countermodel pairs (\cref{sec:pairs}), the eight method obstructions (\cref{sec:methods}), and the arithmetization in \cref{thm:rosser}.

In the manuscript directory (which postdates the \texttt{e4809cf} Lean baseline), \nolinkurl{VERIFICATION.md} gives the theorem-to-declaration map, the checkpoint replay and axiom record (verified against Lean 4.36.0 with only \texttt{propext}, \texttt{Classical.choice} and \texttt{Quot.sound}), and reproduction commands; \nolinkurl{Check.lean} typechecks the stated specializations and axiom dependencies for 30 declarations (explicitly importing \nolinkurl{PartialUniformity}). \nolinkurl{LEDGER.md} documents the literature comparisons and open review questions, and \nolinkurl{CHANGES.md} maps this revision to the archived version~1.1 report (where an endpoint-path extension for overlapping partial-sum intervals and a worked first-precision two-prime \(\{2,5,13\}\) certificate remain). The \repo{scripts/phase17_limitations.py}{countermodel script} checks finite digit, residue, edit, and switching demonstrations for \(C\), \(F\), and a Stoneham oracle (not the normal Liouville stream of \cref{prop:endpoints}). The archived \(\{2,5,13\}\) example is reproducible by the \repo{scripts/overlap_certificate.py}{overlap script}; its finite arithmetic and the analytic density passage have different verification scopes.

\section{Use of AI tools}
This manuscript was produced in an AI-assisted research project. Large language models proposed conjectures and proof strategies, wrote the proofs and the text, implemented the exact computations and Lean 4 formalizations, searched the literature and reviewed drafts across several model providers. The repository owner directed the project and is responsible for its release. The Lean kernel checked the declarations listed in the verification map. The written proofs, literature comparisons and priority statements have been checked only by AI models from different providers. They have not been checked by an independent human expert. Cross-model review is not peer review.

\begingroup\small
\bibliographystyle{plainurl}
\bibliography{references}
\endgroup
\end{document}
